% PNP RESEARCH REPORT::CLAUDE BRIDGE-EQUIV-0001::2026-11-13T01:48:00-08:00
% Independent red-team audit of the shared worst-case derivation-cost directive.
\documentclass[11pt]{article}
\usepackage[T1]{fontenc}
\usepackage[utf8]{inputenc}
\usepackage[margin=1in]{geometry}
\usepackage{amsmath,amssymb,amsthm,booktabs,longtable,array,xurl}
\usepackage[hidelinks]{hyperref}
\setlength{\emergencystretch}{3em}
\newtheorem{theorem}{Theorem}
\newtheorem{proposition}[theorem]{Proposition}
\newtheorem{corollary}[theorem]{Corollary}
\newtheorem{lemma}[theorem]{Lemma}
\theoremstyle{definition}
\newtheorem{definition}[theorem]{Definition}
\newcommand{\poly}{\operatorname{poly}}
\newcommand{\cost}{\operatorname{cost}}
\newcommand{\UNSAT}{\mathrm{UNSAT}}
\newcommand{\Wbar}{\overline W}
\title{The Conditional Bridge $\mathrm{SAT}\in\mathrm P\Rightarrow\Wbar_R(N)\le\poly(N)$\\
is a Reformulation of $\mathrm P\neq\mathrm{NP}$ for Every System\\
Where the Lower-Bound Half Is Available}
\author{Independent red-team run (Claude, external to LSPV-1..5)\\
P vs NP Research Project, repository \texttt{j15marti/pnp}}
\date{13 November 2026, 01:48 PST}
\begin{document}
\maketitle

\section*{Status}
\textbf{P versus NP remains OPEN.} Nothing below is a separation or a step
toward one. This run is an audit of the \emph{shared worst-case derivation-cost
directive} issued to all five roles. Its two main statements
(the bridge-collapse theorem~\ref{thm:collapse} and the
free-strategy theorem~\ref{thm:buys}) are \textbf{PROVED/EXACT} but
elementary: they are bookkeeping about what the directive's own target
\emph{means}, not new mathematics. Their content is negative. The finite
tables in \S\ref{sec:comp} are \textbf{COMPUTATIONAL}.

\medskip\noindent\textbf{One-line finding.} For every fixed proof system $R$
for which an unconditional superpolynomial lower bound on
$\Wbar_R$ is known today, the prioritized conditional bridge
\[
  \mathrm{Bridge}(R):\qquad
  \mathrm{SAT}\in\mathrm P\ \Longrightarrow\ \Wbar_R(N)\le\poly(N)
\]
is \emph{logically equivalent to $\mathrm P\neq\mathrm{NP}$}; and for every
fixed $R$, the conjunction of the directive's two priorities is equivalent to
$\mathrm P\neq\mathrm{NP}$. So the bridge is the weakest \emph{sufficient}
condition, as the directive says, while being no weaker than the goal itself
wherever the other half is deliverable. By the project's own world-ledger
convention (row L13 note in \texttt{AGENTS.md}), a statement equivalent to
$\mathrm P\neq\mathrm{NP}$ is a reformulation, not progress.

\section{Setup and exact quantifiers}

\begin{definition}[Cook--Reckhow system, cost, worst-case free-strategy minimum]
\label{def:cr}
A \emph{proof system for $\UNSAT$} is a relation $V_R(F,\pi)\in\{0,1\}$ that is
(i) \emph{sound}: $V_R(F,\pi)=1$ for some $\pi$ only if the CNF $F$ is
unsatisfiable; (ii) \emph{complete}: every unsatisfiable $F$ has some $\pi$
with $V_R(F,\pi)=1$. It is \emph{Cook--Reckhow} if additionally $V_R$ is
computable in time $\poly(|F|+|\pi|)$. A \emph{cost measure} is any
$\cost_R(\pi)\in\mathbb N$. Write $|\langle F\rangle|$ for the encoded input
length and
\[
 \Wbar_R(N)\;=\;\max_{\substack{F\in\UNSAT\\ |\langle F\rangle|\le N}}
 \ \min_{\pi\,:\,V_R(F,\pi)=1}\ \cost_R(\pi).
\]
$R$ is \emph{polynomially bounded} (p-bounded) iff $\Wbar_R(N)\le\poly(N)$ for
$\cost_R(\pi)=|\pi|$.
\end{definition}

The inner $\min$ is the directive's deliberate design choice: the optimal
derivation sequence is granted free, with no charge for finding it. We adopt
that convention throughout and confirm in Theorem~\ref{thm:buys} that it is a
legitimate lower-bound convention --- and that it is exactly the wrong half to
give away.

\begin{definition}[the prioritized bridge]
For a fixed system $R$ and cost measure, $\mathrm{Bridge}(R)$ is the
$\Pi$-statement
$\mathrm{SAT}\in\mathrm P\Rightarrow \exists c\ \forall N\ \Wbar_R(N)\le N^c+c$.
\end{definition}

Model matching, for the record: $R$ is \emph{fixed in advance}, independent of
the hypothetical polynomial-time SAT algorithm $A$. This is forced by the
directive's own phrasing (``a fixed, sufficiently powerful logical system''),
and it is the hypothesis that drives everything below. If $R$ is allowed to
depend on $A$ --- e.g.\ the system whose proofs are accepting run traces of $A$
--- then $\Wbar_R\le\poly$ is trivial and carries no information.

\section{The bridge collapses onto the goal}

\begin{theorem}[Bridge collapse]\label{thm:collapse}
Let $R$ be any proof system for $\UNSAT$ with $\cost_R(\pi)=|\pi|$.
\begin{enumerate}
\item[\textup{(a)}] If $R$ is Cook--Reckhow and $\Wbar_R(N)\le\poly(N)$, then
$\mathrm{NP}=\mathrm{coNP}$.
\item[\textup{(b)}] $\mathrm P\neq\mathrm{NP}\Rightarrow\mathrm{Bridge}(R)$,
for every $R$, vacuously.
\item[\textup{(c)}] If $\Wbar_R(N)\neq O(N^{c})$ for every $c$ --- i.e.\ the
directive's priority~1 holds for $R$ --- then
$\mathrm{Bridge}(R)\Rightarrow\mathrm P\neq\mathrm{NP}$.
\item[\textup{(d)}] Consequently, if the superpolynomial lower bound of
priority~1 is a \emph{theorem} for $R$, then $\mathrm{Bridge}(R)$ and
$\mathrm P\neq\mathrm{NP}$ are equivalent statements; in particular
$\mathrm{Bridge}(R)$ is \emph{false} in Algorithmica.
\end{enumerate}
\end{theorem}

\begin{proof}
(a) Proofs of size $\le\poly(N)$ verifiable in polynomial time put
$\UNSAT$-CNF, a coNP-complete language, in NP; this is Cook--Reckhow
\cite{cookreckhow79}. (b) $\mathrm{SAT}\in\mathrm P$ iff
$\mathrm P=\mathrm{NP}$, so under $\mathrm P\neq\mathrm{NP}$ the antecedent is
false. (c) Assume $\mathrm{Bridge}(R)$. The consequent fails by hypothesis, so
the antecedent fails: $\mathrm{SAT}\notin\mathrm P$, i.e.\
$\mathrm P\neq\mathrm{NP}$. (d) Combine (b) and (c); falsity in Algorithmica is
(c) read contrapositively under $\mathrm P=\mathrm{NP}$.
\end{proof}

\begin{corollary}[no decomposition gain]\label{cor:nogain}
For every fixed $R$, the conjunction
\[
 \bigl[\ \Wbar_R(N)\ \text{superpolynomial}\ \bigr]\ \wedge\ \mathrm{Bridge}(R)
\]
is equivalent to $\mathrm P\neq\mathrm{NP}$: it implies it by
Theorem~\ref{thm:collapse}(c), and is implied by it whenever the first
conjunct is a theorem. Hence splitting the goal into ``lower bound'' plus
``weakest sufficient bridge'' transfers difficulty rather than reducing it:
all systems with a known lower bound put \emph{all} remaining difficulty into
a statement that is a verbatim restatement of the goal.
\end{corollary}

\noindent The hypothesis of (d) is satisfied by every system for which the
project could realistically discharge priority~1 today:

\begin{center}
\begin{tabular}{@{}lll@{}}
\toprule
System $R$ & unconditional superpoly.\ lower bound & status of $\mathrm{Bridge}(R)$\\
\midrule
Resolution & $2^{\Omega(n)}$, $\mathrm{PHP}^{n+1}_n$ \cite{haken85} & $\equiv \mathrm P\neq\mathrm{NP}$\\
Tree-like resolution & $n^{\Omega(n)}$, $\mathrm{PHP}$ \cite{dantchevriis01} & $\equiv \mathrm P\neq\mathrm{NP}$\\
Res$(k)$, $k$ fixed & known & $\equiv \mathrm P\neq\mathrm{NP}$\\
Nullstellensatz, PC, PCR (fixed field) & known & $\equiv \mathrm P\neq\mathrm{NP}$\\
Cutting planes & $2^{n^{\varepsilon}}$ \cite{pudlak97} & $\equiv \mathrm P\neq\mathrm{NP}$\\
Bounded-depth Frege & $2^{n^{\varepsilon}}$ \cite{ajtai94,pbi93} & $\equiv \mathrm P\neq\mathrm{NP}$\\
\midrule
$\mathrm{TC}^0$-Frege, Frege, Extended Frege & \textbf{OPEN} & \textbf{OPEN}, not known equivalent\\
\bottomrule
\end{tabular}
\end{center}

For the last line the equivalence argument does not apply, because the
consequent $\Wbar_R\le\poly$ may be true; this is exactly the regime where the
decomposition carries information, and exactly the regime where priority~1 is
the open frontier of proof complexity. We record the resulting dilemma:

\begin{corollary}[strength dilemma]\label{cor:dilemma}
For every fixed Cook--Reckhow $R$, exactly one of:
\begin{itemize}
\item $R$ is p-bounded: then $\mathrm{NP}=\mathrm{coNP}$ and priority~1 is
\emph{impossible} for $R$ (no superpolynomial bound exists); or
\item $R$ is not p-bounded: then $\mathrm{Bridge}(R)$ implies
$\mathrm P\neq\mathrm{NP}$, so no proof of $\mathrm{Bridge}(R)$ can be easier
than a proof of the goal.
\end{itemize}
In neither case can both halves be ``cheap''. The directive's hope that the
bridge is the easy half is therefore unavailable for \emph{every} system, not
merely for weak ones --- what varies with strength is only \emph{which} half is
hard.
\end{corollary}

\section{What $\mathrm P=\mathrm{NP}$ actually buys is the part the
directive gives away}

\begin{theorem}[free strategy is the only known consequence]\label{thm:buys}
Let $R$ be Cook--Reckhow and suppose $\mathrm{SAT}\in\mathrm P$. Then there is
a deterministic algorithm that, on unsatisfiable input $F$, outputs a
\emph{minimum-length} $R$-refutation of $F$ in time
$\poly\bigl(|\langle F\rangle| + \min_{\pi}|\pi|\bigr)$. By contrast, no
bound whatsoever on $\min_\pi|\pi|$ follows from $\mathrm{SAT}\in\mathrm P$ for
a fixed $R$; the existence of \emph{some} p-bounded system does follow
(Cook--Reckhow, via $\mathrm{NP}=\mathrm{coNP}$), but that system is not $R$.
\end{theorem}

\begin{proof}
$L_R=\{(F,1^s):\exists\pi,\ |\pi|\le s,\ V_R(F,\pi)=1\}$ is in NP since $V_R$
is polynomial-time. Under $\mathrm{SAT}\in\mathrm P$ we have
$\mathrm P=\mathrm{NP}$, so $L_R\in\mathrm P$; doubling $s$ locates the
minimum length $s^\ast$ with $O(\log s^\ast)$ queries, and prefix
self-reduction on the witness (also in P under $\mathrm P=\mathrm{NP}$)
extracts $\pi$ symbol by symbol in $\poly(|F|+s^\ast)$ steps. The second
assertion is the content of Theorem~\ref{thm:collapse}(d): for $R$ with a
known lower bound, $\mathrm{SAT}\in\mathrm P$ \emph{contradicts}
$\Wbar_R\le\poly$ rather than implying it.
\end{proof}

\noindent\textbf{Red-team reading.} The directive states, correctly, that
``finding the optimal sequence can only add work, so we do not need to charge
for that search''. That is sound as a lower-bound convention. But
Theorem~\ref{thm:buys} says that \emph{proof search} is precisely and only
what a polynomial-time SAT algorithm is known to deliver for a fixed proof
system: $\mathrm P=\mathrm{NP}$ makes every Cook--Reckhow system
automatable, and nothing more. By discarding the search charge, the program
discards the single known link between its hypothesis and its object, and
retains an obligation (a size bound for a fixed $R$) that is equivalent to
$\mathrm{NP}=\mathrm{coNP}$ in the existential form and open in the fixed-$R$
form. The asymmetry is not fixable by choosing a cleverer $R$: the one known
unconditional-looking hardness statement about proof search,
non-automatability of resolution, is itself only \emph{NP}-hardness
\cite{atseriasmuller20}, hence consistent with $\mathrm P=\mathrm{NP}$, and an
unconditional non-automatability theorem would already imply
$\mathrm P\neq\mathrm{NP}$.

\section{Which cost measures can carry the target at all}
\label{sec:cost}

The directive assigns the Time--Space Resource Auditor ``inference costs, proof
encodings, computational overhead, and peak memory''. Peak memory cannot carry
the target:

\begin{proposition}[cost-measure collapse]\label{prop:collapse}
If a cost measure satisfies
$\min_{\pi}\cost_R(\pi)\le\poly(|\langle F\rangle|)$ for all unsatisfiable $F$,
then $\Wbar_R(N)\le\poly(N)$ unconditionally and priority~1 is impossible for
that measure. This holds for: resolution \emph{clause space}
($\le n+2$ on $n$ variables, via the DPLL tree), resolution \emph{width}
($\le n$), derivation \emph{depth} ($\le n$), and any measure dominated by
these. It fails, and only fails, for length-like (size, weight, total-symbol)
measures.
\end{proposition}

\begin{proof}
Immediate from the definition of $\Wbar_R$ as a minimum over all $\pi$: a
single polynomial-cost refutation per instance caps the minimum. The clause
space bound is the standard DPLL simulation; depth $\le n$ because a complete
decision tree over the $n$ variables falsifies some clause on every branch.
\end{proof}

Note the matching known lower bound: Tor\'an's $\Omega(n)$ clause-space bound
for $\mathrm{PHP}$ \cite{toran99} is already optimal against the trivial $n+2$
ceiling, so there is no room in a space measure. Any ``minimal work to reach
the minimal reduced logical state'' intended as peak memory is therefore
$\poly(N)$ by construction, and the project's multisign-DAG objective must be
read as a length-like measure to be non-vacuous.

\section{Computational check: free-strategy minima are small and well-defined}
\label{sec:comp}

Two independent stages, each well under the 30-second project runtime rule,
compute \emph{exact} free-strategy minima (minimum over \emph{all} refutations,
no charge for search) for the canonical hard families at sizes where
exhaustive minimization is feasible. Programs:
\texttt{experiments/claude\_wr/exact\_treeres\_min.py} and
\texttt{experiments/claude\_wr/cost\_measure\_collapse.py}; raw data:
\texttt{results/raw/claude\_wr/2026-11-13\_0148\_treeres\_exact.json},
\texttt{results/raw/claude\_wr/2026-11-13\_0148\_cost\_measure\_collapse.json}.
Chunking strategy: one instance per stage invocation, memoized recursion over
partial assignments, instances ordered by variable count so that each
individual call terminates in seconds.

\begin{center}
\begin{tabular}{@{}lrrrrrr@{}}
\toprule
Instance & vars & clauses & $N$ & min leaves & min nodes & min depth\\
\midrule
$\mathrm{PHP}^{2}_{1}$ & 2 & 3 & 7 & 3 & 5 & 2\\
$\mathrm{PHP}^{3}_{2}$ & 6 & 9 & 27 & 11 & 21 & 4\\
Tseitin $K_4$ (odd charge) & 6 & 16 & 64 & 28 & 55 & 5\\
$\mathrm{PHP}^{4}_{3}$ & 12 & 22 & 70 & 43 & 85 & 7\\
\bottomrule
\end{tabular}
\end{center}

Here ``min leaves'' is the exact minimum number of leaves of a decision tree
solving the false-clause search problem, i.e.\ the exact minimum tree-like
resolution refutation leaf count; ``min nodes'' $=2\cdot\text{leaves}-1$;
``min depth'' is the exact minimum tree depth. Readings, stated as negatives:

\begin{itemize}
\item The inner minimum of the directive's objective is a well-defined,
computable quantity at small scale --- the objective is not ill-posed.
\item These numbers are \emph{not} evidence of superpolynomial growth; they are
far below $N$ at every row and are consistent with the known asymptotics
($n^{\Omega(n)}$ for tree-like PHP \cite{dantchevriis01}). Finite tables of
this kind cannot distinguish the hypotheses at issue, and no run should cite
them as support.
\item Min depth $\le$ \#vars in every row, as Proposition~\ref{prop:collapse}
requires; clause space $\le$ depth $+2\le 9$ even where the minimum tree has
43 leaves. The size/space gap is visible already at $N=70$.
\end{itemize}

\section{Literature check for the current obstacle}

Searched this run (primary sources; snippets only, full texts not re-verified
here):
\begin{itemize}
\item Cook--Reckhow \cite{cookreckhow79}: a p-bounded system exists iff
$\mathrm{NP}=\mathrm{coNP}$. This is the exact wall in front of priority~1+3 as
stated, and the reason Theorem~\ref{thm:collapse}(a) is unavoidable.
\item Kraj\'i\v{c}ek--Pud\'l\'ak and successors: $\mathrm{NE}=\mathrm{coNE}$
implies optimal proof systems, $\mathrm E=\mathrm{NE}$ implies p-optimal ones;
weakened to (N)EE by K\"obler--Messner--Tor\'an. Existence of (p-)optimal
systems for TAUT remains \textbf{OPEN}, and Egidy--Gla\ss er (STOC 2025) plus
Egidy (arXiv:2602.02294, preprint, unverified) construct oracles showing there
is no relativizing positive proof. \emph{Relevance:} any attempt to upgrade
``some p-bounded system'' to ``our fixed $R$ is p-bounded'' runs into this
open problem, and a relativizing argument cannot do it.
\item Kraj\'i\v{c}ek (ECCC TR21-053): a time-optimal proof search algorithm
with respect to all systems exists iff a p-optimal system exists. \emph{This is
the precise home of the directive's ``free optimal sequence'' grant.}
\item Atserias--M\"uller \cite{atseriasmuller20}: resolution is not
automatable unless $\mathrm P=\mathrm{NP}$. Used in
Theorem~\ref{thm:buys}'s discussion; it is a conditional statement and
therefore cannot be combined with $\mathrm{SAT}\in\mathrm P$ to a
contradiction.
\item Pich--Santhanam (arXiv:2312.08163): a conditional route to
$\mathrm P\neq\mathrm{NP}$ from the unprovability of circuit lower bounds in
Extended Frege, assuming among others that $S^1_2$ proves a polynomial-time
function generates anticheckers for SAT. \emph{Relevance:} this is the live
literature on precisely the top row of our table --- the only regime where
$\mathrm{Bridge}(R)$ is not a restatement of the goal --- and it proceeds by
bounded-arithmetic provability of soundness/reflection, not by derivation-cost
minimization.
\end{itemize}

\textbf{Transfer classification.} (i) \emph{Directly transferable theorem:}
Cook--Reckhow \cite{cookreckhow79}, used as stated.
(ii) \emph{Useful structural analogy needing a new transfer theorem:} the
bounded-arithmetic soundness route (Kraj\'i\v{c}ek--Pud\'l\'ak simulation
criterion, Pich--Santhanam) --- it would have to be connected to a
multisign-DAG cost measure by an explicit object correspondence that does not
exist yet. (iii) \emph{Misleading analogy / failed transfer:} reading
``minimal work to reach the minimal reduced logical state'' as peak memory
(Proposition~\ref{prop:collapse} kills it), and reading NP-hardness of
automating resolution as a step toward an unconditional lower bound
(Theorem~\ref{thm:buys}).

\textbf{Post-quantum technique ledger (required consideration).} Items 1--4
(Schur squares, quadratic relation spaces, Hilbert functions, collision
entropy) act on \emph{algebraic} encodings and are orthogonal to
Theorem~\ref{thm:collapse}, which is independent of $R$'s internal structure:
no choice of algebraic invariant changes the fact that a fixed Cook--Reckhow
system either is or is not p-bounded. Item~5 (worst-case/average-case
architecture, LWE lesson) is the one with live content here, and
Theorem~\ref{thm:buys} sharpens it: external extraction from black-box solver
behaviour is the right shape, but what is extractable under
$\mathrm{SAT}\in\mathrm P$ is a \emph{search} procedure, not a size bound.
Classification: items~1--4 \textbf{FAILED TRANSFER} for this obstacle,
item~5 \textbf{ANALOGY/CANDIDATE}, items~6--7 \textbf{FAILED TRANSFER}
(assumption-bearing).

\section{Corpus reviewed}

Repository head \texttt{6a9ccb7c} on \texttt{main}
(2026-11-13 01:45:43 -0800), working tree clean. Surveyed: root
\texttt{AGENTS.md} (all allocation headings, including the 2026-11-09
five-role knot/link topology handoff and the 2026-10-28 world-ledger pointer),
root \texttt{main.tex} ownership-marker structure (25{,}750 lines; LSPV-1..5
regions located, none edited by this run), \texttt{tasks/} specifications
(TASK-001..007), the coordination handoff index
\texttt{research/coordination/} (newest: LSPV-3/609 native-accept goal
separation, 01:24 PST), the newest peer artifacts of each role
(LSPV-1/613 charge congruence, LSPV-2/594 residual LP, LSPV-3/609 dual-goal
no-go, LSPV-4/581 witness-work firewall, LSPV-5/521 control-namespace
red audit), \texttt{research/agents/codex/} reports of 2026-11-11 (the
only prior artifacts in the corpus using the $\Wbar$/$H$ minimum-work
objective), and the world ledger
\texttt{research/agents/claude/world-ledger/world\_ledger.csv} (rows L01--L40).

\textbf{Duplication / contradiction audit.} The $\Wbar$-style objective
appears previously in \texttt{research/agents/codex/2026-11-11\_1203} and
\texttt{\_1331}, and in the $H(F)$ benchmark of
\texttt{2026-11-11\_2035\_COMPATIBILITY\_CORE\_AND\_DEGREE\_TRANSFER.tex},
whose status line already records that ``the all-program work amplification
needed for a separation remains OPEN''. No contradiction with the present
report: those runs bound $\overline H$ \emph{above} from one fixed decider,
which is the direction Theorem~\ref{thm:collapse} leaves untouched. This run
adds the missing statement about the \emph{conditional} direction. No
duplicate of Theorem~\ref{thm:collapse} was found in the corpus by search for
\texttt{overline W\_R}, \texttt{W\_R(N)}, or bridge/transfer phrasing.

\textbf{What changed relative to root \texttt{main.tex}.} Nothing. This run is
external to the five LSPV roles and edits no ownership region; per
\texttt{AGENTS.md}, structural reconciliation of root \texttt{main.tex} is
LSPV-5's. The finding is offered to LSPV-5 and to the General SAT Lower-Bound
Lead as an input, with the recommendation in \S\ref{sec:rec}.

\textbf{Adopted / rejected / retired prior components.} Adopted: the project's
Cook--Reckhow discipline and status vocabulary; the codex $H(F)$ upper-bound
form as the complementary direction. Rejected for this obstacle: space/width
/depth cost measures (Proposition~\ref{prop:collapse}). Retired: nothing ---
this run retires no prior result.

\textbf{Stale or overstated claims found.} One, in the directive itself rather
than in a report: the characterization of
$\mathrm{SAT}\in\mathrm P\Rightarrow\Wbar_R\le\poly(N)$ as a ``weakest
sufficient bridge'' that may be ``proved directly''. It is weakest among
sufficient conditions and simultaneously equivalent to the goal for every
system with a known lower bound (Theorem~\ref{thm:collapse}(d)); it admits no
direct proof that is not a proof of $\mathrm P\neq\mathrm{NP}$. No agent report
in the corpus asserts the bridge as established, so no published status label
requires downgrading.

\section{Recommendation}\label{sec:rec}

\begin{enumerate}
\item Keep the objective $\Wbar_R$, but state alongside it which of the two
regimes a run is in. Below the frontier (resolution through bounded-depth
Frege) a run must not present bridge work as progress: it is
$\mathrm P\neq\mathrm{NP}$ in other words. Above the frontier (Frege, EF,
$\mathrm{TC}^0$-Frege, unrestricted-sharing multisign DAGs) the lower-bound
half is the open problem and no finite hard family settles it.
\item Fix the cost measure to a length-like one explicitly, and drop peak
memory from the shared objective (keep it as an auditing quantity only).
\item If the bridge is to be attacked above the frontier, the live technique in
the literature is bounded-arithmetic provability of soundness/reflection, not
derivation-cost minimization; the missing object correspondence between a
multisign-DAG cost and an EF-provable reflection principle is the thing to
write down.
\item Ledger: record the bridge as an \emph{equivalent of}
$\mathrm P\neq\mathrm{NP}$ (rows L41--L44 added this run), so that future runs
classify it correctly on sight.
\end{enumerate}

\begin{thebibliography}{9}
\bibitem{cookreckhow79} S.~A. Cook and R.~A. Reckhow.
\emph{The relative efficiency of propositional proof systems.}
J.\ Symbolic Logic 44(1):36--50, 1979.
\bibitem{haken85} A.~Haken. \emph{The intractability of resolution.}
Theoretical Computer Science 39:297--308, 1985.
\bibitem{dantchevriis01} S.~Dantchev and S.~Riis.
\emph{Tree resolution proofs of the weak pigeon-hole principle.}
Proc.\ CCC, 2001.
\bibitem{pudlak97} P.~Pud\'l\'ak. \emph{Lower bounds for resolution and
cutting plane proofs and monotone computations.}
J.\ Symbolic Logic 62(3):981--998, 1997.
\bibitem{ajtai94} M.~Ajtai. \emph{The complexity of the pigeonhole principle.}
Combinatorica 14(4):417--433, 1994.
\bibitem{pbi93} T.~Pitassi, P.~Beame, R.~Impagliazzo.
\emph{Exponential lower bounds for the pigeonhole principle.}
Computational Complexity 3:97--140, 1993.
\bibitem{toran99} J.~Tor\'an. \emph{Lower bounds for space in resolution.}
Proc.\ CSL, 1999.
\bibitem{atseriasmuller20} A.~Atserias and M.~M\"uller.
\emph{Automating resolution is NP-hard.} J.\ ACM 67(5):31, 2020.
\bibitem{krajicek21} J.~Kraj\'i\v{c}ek. \emph{Information in propositional
proofs and algorithmic proof search.} ECCC TR21-053, 2021.
\bibitem{pichsanthanam23} J.~Pich and R.~Santhanam.
\emph{Towards P $\neq$ NP from extended Frege lower bounds.}
arXiv:2312.08163, 2023.
\end{thebibliography}

\medskip
\noindent\textbf{P versus NP remains OPEN.} The results of this run are
negative and definitional; they close no proof arrow and open no new one.

% END_OF_VERIFIED_REPORT::CLAUDE::BRIDGE-EQUIV-0001
\end{document}
